% Provenance: copied from the design workspace CambridgePQ/recursion.fast, `report/sections/06-architecture.tex`, on 2026-10-19. Read-only copy; the source is authoritative.
% 06-architecture -- written by W2
\section{Competition architecture and trust model}\label{sec:architecture}

This section describes how a submission travels from an untrusted pull request to a record, which
software and which mathematics a reader must trust along the way, and the budgets and versioning
rules that keep the process reproducible. The machinery is inherited from ots.golf and sig.golf
(\cref{sec:prior}); the new elements are the rendering of raw RISC-V images and of every score
input into the challenge statement, a hygiene and computability gate that inspects the elaborated
Lean environment, score functions that are Lean definitions of the contract, and an unscored badge
path that runs candidate binaries.

\subsection{Repositories and roles}\label{subsec:arch-repos}

\paragraph{Two repositories.} The \emph{contract repository} (\repo{recursion.fast}) is trusted. It
holds the Lean contract, the challenge stubs, the harness and verifier scripts, the pinned
toolchain manifest and the controls (\cref{subsec:sec-nonvacuity}). It contains no submission
proofs. The organizer library (machine facts, symbolic executor, field routines, IOP components,
bridges to ArkLib and CompPoly) lives in the same repository as a separate package that is
\emph{not} protected: it is kernel-checked like any submission and may be imported by agents, but
nothing in it is trusted. The \emph{submissions repository} (\repo{recursion.fast-submissions}) is
untrusted. Its \code{main} branch holds the current record root per track, \code{records.json},
contributor metadata and a pin of the contract commit for local checking.

\paragraph{Intake is pull-request only, and nothing is merged.} A pull request may create or change
exactly one submission root. The verification bot fetches the head commit at a frozen hash, retains
it under a creation-only tag, keeps only that root and verifies it against the trusted contract
checkout. Pull requests are never merged or closed by the bot. After a durable verdict the bot
copies the checked root and its registry entry into \code{main} in a separate commit. Receipt and
verdict comments are the durable record; the verification server is disposable and can be rebuilt
from the repositories. This is ots.golf's protocol \cite{otsgolf2026}. It answers two attacks: a
merge would let a submission modify files outside its root (including scripts that later runs
execute), and a mutable head would let the verified commit differ from the published one.

\paragraph{Record rule.} A record is the first verified head whose certified score strictly
improves the current record of its track and board (unconditional or conditional, per class
display) at the time its verification \emph{finishes}, and which passes the badge path
(\cref{subsec:arch-badge}). A later identical or copied claim never takes a record. Every root
carries a \code{NOTES.md} (idea, result, what failed, what to try next), published as untrusted
text whatever the verdict, so that the autoresearch loop shares negative results.

\paragraph{Roles.}
\begin{description}[nosep,leftmargin=1.2em]
\item[Organizers] maintain the contract, the harness and the organizer library; publish versions;
  port records across contract changes. Changes to the contract are reviewed by two organizers.
\item[Hypothesis reviewers] (at least two, one a coding theorist) admit, version and retire named
  mathematical hypotheses for the conditional track (\cref{subsec:sec-hypotheses}).
\item[Submitters] (human teams or autonomous agents) submit roots; they control everything inside
  their root and nothing else.
\item[Auditors] read the contract, the trusted-surface list and the controls; their review is the
  only defense against statement-level holes, which no checker detects (\cref{subsec:bg-lean}).
\item[Harness operators] run the verification and badge services on pinned hardware.
\end{description}

\subsection{The submission bundle}\label{subsec:arch-bundle}

\begin{table}[t]
\centering\small
\begin{tabularx}{\linewidth}{@{}>{\raggedright\arraybackslash}p{3.3cm}>{\raggedright\arraybackslash}X>{\raggedright\arraybackslash}p{3.0cm}@{}}
\toprule
File(s) & Content and role & Trust \\
\midrule
\code{images/}\allowbreak\code{verifier.code}, \code{images/}\allowbreak\code{verifier.data} & raw little-endian 32-bit instruction words and 64-bit data words, plus a \code{harts} flag; T2 has embed, fold and decide images & untrusted; rendered into the challenge, so the theorems are about these bytes \\
\code{claim.json} & per class the 18 claimed counters $c_k$ (T2: three vectors), the declared shape summary (rounds, message rounds, $E$-challenges, index challenges, total message length, tape input words), the claimed scores; latency board: $(W_k,S_k)$ & untrusted; strict typed schema; every value is rendered into the statement and proved \\
\code{Solution.lean}, \code{RecursionFast}\allowbreak\code{Candidate/**} & the \lean{submission} (protocol per class: indexer, honest prover, verifier, extractor, errors) and the proof of the certificate & checked by the gate, comparator and kernel replays \\
native indexer and prover & Rust source built by the harness, or a static x86-64 binary & untrusted, sandboxed; badge path only \\
\code{NOTES.md} & research notes & untrusted text, published \\
\bottomrule
\end{tabularx}
\caption{Contents of a submission root.}
\label{tab:arch-bundle}
\end{table}

\Cref{tab:arch-bundle} lists the files. Two choices need justification. First, the verifier is
submitted as \emph{raw bytes} rather than as a Lean term: the trusted emulator of the badge path and
any downstream deployment must run exactly the program the theorem is about, and sig.golf found that
extracting images from candidate Lean by \code{\#eval} runs code the candidate controls
(\lean{implemented\_by}, evaluation overrides, files written at build time), so it removed the
extraction \cite{siggolf2026}. ELF files are rejected because an ELF parser would join the trusted
base; a trusted \code{elf2image} script will be offered to agents instead. Second, the honest prover is a
Lean \emph{specification}, used only for completeness. A RISC-V prover with a refinement proof would
be the largest proof in any submission; prover efficiency is gated practically by the native prover
on the badge path.

\subsection{The binding rule}\label{subsec:arch-binding}

\begin{decision}{A-1}{Bytes and score inputs reach the theorem as organiser data}
The trusted renderer reads \code{images/*.code} and \code{images/*.data} with size checks and
emits them into the generated challenge as chunks of 256 words, each packed into one natural-number
literal and decoded by the contract's \lean{unpack} into an \lean{Image}
(\lean{List (BitVec 32)}, \lean{List Word}, \lean{harts}). The image, the declared shape summaries,
the claimed counters and the claimed scores are fields of the rendered \lean{Declared1} (T0:
\lean{Declared0}, T2: \lean{Declared2}); the solution provides only \lean{submission}, and the
certificate's \lean{summary} and \lean{score} clauses make it prove that its protocols have the
rendered summaries and that the rendered scores are the contract's score functions of the rendered
data. The badge path's loader reads the \emph{same files}. Hence: bytes $\to$ renderer $\to$ Lean
literal $\to$ theorem, and the same bytes $\to$ loader $\to$ emulator; nothing is evaluated from the
solution.
\end{decision}

\noindent This is the v0.3 form of the rule (A4-F2). In v0.2 the image was a field of the submission
and the shape inputs of the score were extracted from the solution, so an \lean{implemented\_by}
override could make compiled evaluation report a different image or round count from the one the
kernel checked (probe \lean{UnboundImage}). The renderer accepts a strict schema (A4-F5): typed
placeholders \code{\{\{nat:k\}\}} (a JSON integer, not a boolean, below $2^{64}$),
\code{\{\{rat:k\}\}} (a string \code{a} or \code{a/b}) and \code{\{\{image:k\}\}} (an object with
exactly \code{code}, \code{data}, \code{harts}, code plus data $\le2^{20}$ bytes); the claim's key
set must equal the stub's placeholder set. The image
paths in a claim are confined (A6-R1): each must be a plain relative path (not absolute, no
\code{..}, \code{.} or empty component, no backslash), free of symbolic links on every component,
and name a regular file whose real path lies inside the claim's real directory. Before this rule the
renderer followed \code{os.path.join(base, path)}, and A6's probe claim rendered a harness secret
from \path{/tmp} into the public challenge text as image words; any harness-readable file of
suitable length, including another team's pending image, could be published this way. The
renderer tests now hold 22 bad claims (13 schema cases, 7 path cases, A6's two probes), all
rejected. Costs are
measured at the 1\,MiB cap ($2^{18}$ words). The chunked challenge elaborates in 13\,s with 2.0\,GB
(a flat list literal took more than 250\,s and 7\,GB). A lookup \lean{image.code[j]? = some w} by
\lean{decide +kernel} walks the list prefix, about 80\,s and 9\,GB \emph{per lookup}, which makes a
proof with one decode fact per instruction quadratic; the library lemma \lean{getElem?\_chunks'}
reduces a lookup to one chunk after a one-time 28\,s length check. \lean{byteSize ≤ imageCap}
still costs about 80\,s and 9\,GB once per certificate.

\subsection{Toolchain pin}\label{subsec:arch-toolchain}

\begin{decision}{A-2}{One toolchain for contract, library and the components agents reuse}
Lean v4.34.1 and Mathlib v4.34.1; ArkLib \texttt{7653a90}, CompPoly \texttt{572f997} and VCVio
\texttt{a4232d0}, which is leanerVM's dependency set \cite{verifiedzkevm}; riscv-zkvm
\code{Rv64/*} vendored byte-identically from \texttt{6810fc1} with a check script. The contract
imports only Lean core, Mathlib and the vendored model; ArkLib and CompPoly appear only in the
untrusted bridges and library.
\end{decision}

\noindent The pin follows the binary-field components. leanerVM's sumcheck, batching, GKR grand
product, query-check and sequential-composition ports, and CompPoly's \lean{BF64} and
\lean{BF64.Ext3}, are built on v4.34.1. ArkLib's legacy security layer at this commit is
byte-identical to the one reviewed at \texttt{db7005f}. riscv-zkvm, natively on v4.33.1, builds
unchanged on v4.34.1. A later bump costs every submission a port: ots.golf changed its pin 46
times in its first ten days. The pin should therefore be frozen before launch.

\subsection{What continuous integration verifies}\label{subsec:arch-ci}

The verifier runs the following steps for every pull request. Steps 1--7 decide the verdict and the
score; step 8 decides record eligibility. Their \emph{order} is part of the trust model
(Decision~A-3). \emph{Implementation status in v0.3:} \code{scripts/render-challenge.py}
implements step~3; \code{scripts/harness.sh} implements steps 4--6 (stages 0--2), the
fresh-module-root check of step~2, and a hook that runs a badge command as stage~3. Steps 1 and 2
otherwise, the badge checks of step~8, the trusted emulator, the Fiat--Shamir prefix and the
challenge sampler are specified here but not yet implemented. \code{scripts/ci.sh} runs the
organiser-side equivalent (source checks, organiser build, manifest
check, axiom audit, score probes, renderer tests, every stub rendered and elaborated, the toy
certificate through the harness positively and negatively, the A4 gate probes and the A6
regressions, a final manifest check) in 614\,s after round 3 and 685\,s after round 4 (the extra time is the new
regression tests; independent re-run 687\,s), of which about 353\,s are the toy runs through the harness
(v0.3 before round~3: 411\,s).
\begin{enumerate}[leftmargin=1.5em,itemsep=2pt]
\item \textbf{Preflight and export} (planned). Refuse to run as root; require the Linux sandbox (systemd
  transient unit, Landlock via landrun, a seccomp filter denying network and debugging system
  calls, private PID, IPC and device namespaces, a bounded work volume, read-only package mounts).
  Export only the submission root from the pinned commit at blob level (regular files only; no
  symlinks, submodules or git attributes).
\item \textbf{Source policy} (planned, except fresh module roots). Bounded file count and sizes; only the files of
  \cref{tab:arch-bundle}; Lean headers may import only the exact contract modules, the organizer
  library, Mathlib and the pinned dependencies of the library, and the submitter's own modules; no
  \lean{prelude} or \lean{module} headers. Exact module names, not prefixes, are allowed for
  third-party code, because ots.golf was exposed to a sorried instance reached through a
  prefix-allowed import \cite{otsgolf2026}. Module roots must be fresh: a root that already exists in
  the organiser package, the packages or the toolchain is refused, since Lean resolves a module
  through the first search-path root that has its name. A comment-stripped text lint remains as
  defense in depth only: string literals containing \code{/-} hide an \lean{implemented\_by} or a
  \lean{sorry} from it (A4-F3), so the binding checks are steps 5 and 6.
\item \textbf{Render.} Generate the challenge module from the stub, the claim and the image files
  (\cref{subsec:arch-binding}). Image paths are confined to the claim's directory
  (\cref{subsec:arch-binding}, A6-R1).
\item \textbf{Integrity of the organiser build} (stage~0). The organiser package (contract, vendored
  \code{Rv64}, library, controls, bridge, \code{gate}, \code{axiomaudit}) is built by the organisers
  into its own directory, outside every tree in which solution code runs, and is mounted read-only.
  Before any solution code runs, every \code{.olean}, \code{.olean.server} and \code{.olean.private}
  of the toolchain, the packages and the organiser package is checked against a committed SHA-256
  manifest: 35,093 entries (7,556 core, 27,508 packages, 29 contract and vendored), 10\,s. The
  check reports changed, missing \emph{and} extra files, so a dropped-in \code{Mathlib/Evil.olean}
  fails it. Lean builds are bit-reproducible here (a clean rebuild in a copy matched all entries),
  so the manifest is a pin: CI fails when the organiser build drifts.
\item \textbf{Comparator and kernel replays, first} (stage~1). In a \emph{fresh} tree that contains
  only the rendered challenge, the solution sources and a generated lakefile that requires the
  organiser package by path, comparator \cite{comparator} builds and exports the challenge with
  \code{lean4export} \cite{lean4export} \emph{before} it builds the solution, then exports the
  solution, checks that the target theorem has the same type as in the challenge and that every
  constant transitively used by the statement is identical, audits the axiom closure, and replays
  the solution into a fresh kernel environment. The challenge export is held in memory, so no later
  build can alter it. A second kernel, nanoda \cite{nanoda}, replays the same export. The sandbox
  makes only the fresh tree's \code{.lake} writable. Afterwards the manifest is re-checked, and the
  solution's build output may not contain any organiser module path. Options such as
  \lean{debug.skipKernelTC} leave no trace in an \code{.olean}, so the gate cannot see them; the
  replays are the binding check. The evidence so far is \code{leanchecker} in CI (probe
  \lean{SkipKernel}: a false theorem passes the gate and is rejected by \code{leanchecker}); the probe
  has not yet been run through comparator and the harness, which is a planned regression. On the toy challenge this stage takes 115\,s with the default
  kernel and 175\,s with nanoda, under the sandbox (v0.3 in a shared tree: 86\,s and 148\,s; the
  difference is Lake's first load of a new workspace, about 20\,s, plus the sandbox). Changing one
  rendered image word is rejected.
\item \textbf{Environment gate and axiom audit} (stage~2, \code{scripts/Gate.lean}, 354 lines, Lean
  core only). In a \emph{second} fresh tree the solution is built in the sandbox, the manifest and the
  shadow check run again, and the gate loads the \emph{elaborated} environment from the \code{.olean}
  files, so no source text is read. It walks the transitive closure of the certificate theorem
  (types, values, constructors, recursors) and rejects if
  \begin{enumerate}[label=(\roman*),nosep]
  \item the closure contains \lean{sorryAx}, \lean{Lean.ofReduceBool}, \lean{Lean.ofReduceNat},
    \lean{Lean.trustCompiler} or an axiom outside \lean{propext}, \lean{Classical.choice},
    \lean{Quot.sound} (\lean{decide +native} appears as an auxiliary axiom);
  \item a closure constant from an untrusted module is \lean{unsafe}, \lean{partial} or
    \lean{opaque};
  \item \emph{any} untrusted module registers an entry of one of the attributes
    \lean{implemented\_by}, \lean{extern}, \lean{csimp}, \lean{init}, \lean{builtin\_init} or
    \lean{export}, whatever its target;
  \item a badge-path field of the named submission has no compiled code: \lean{VK}, \lean{IDX},
    the prover and the extractor of each protocol (for T2, each step's prover and extractor). Every
    other field (\lean{ε}, \lean{verify}, shapes, proofs) may be \lean{noncomputable}.
  \end{enumerate}
  Rule (iv) guarantees only that compiled code \emph{exists}. In Lean, \lean{noncomputable} marks
  definitions without compiled code (uses of choice, real inverses), not expensive ones: a
  satisfiability table for the index oracle is computable by enumeration over the finite witness
  space and passes (iv). Feasibility of the indexer (A2-M4) is enforced only by the badge path, which
  runs it under a time and memory budget and gates records, not the proved score
  (\cref{subsec:proto-indexer}). Two round-4 fixes (A7, NOTICE-P2e) close gaps that did not affect
  the proved score. Rule (iv) used to check all of \lean{submission}, so an error written as
  \lean{errorCap} or with \lean{ℝ≥0} inverses failed it (A7-M2); it now checks only the badge-path
  fields. And stage~2 rebuilds the solution from source, which runs solution code again; the harness
  now records the hashes of every \code{.olean}, \code{.ir} and \code{.c} artefact that comparator
  verified and rejects a gate-tree build that is not byte-identical (A7-M3, regression
  \code{bindproj}), so the gate and the badge see the verified build. Elaboration must therefore be
  deterministic and independent of the build directory. The \lean{init} and \lean{export}
  rules of (iii) are new in round~3: \lean{@[init]} on an ordinary \lean{def} makes compiled code
  read a value the kernel never sees, and \lean{@[export]} allows symbol interposition in a linked
  badge executable (A6-G1). A solution may therefore not use \lean{initialize},
  \lean{register\_simp\_attr} or \lean{@[export]}. A module is \emph{trusted iff every one of its
  \code{.olean} parts hashes to its manifest entry}; a module whose path is in the manifest but whose
  files differ is rejected as tampered. Trust is neither by name nor by location: a name-prefix rule
  let a solution module named \lean{RecursionFast.Contract.Spoof} pass (probe \lean{SpoofCert}), and
  the location and import-graph rule that replaced it trusted files the solution build could rewrite
  (A6-G2). On the toy certificate the gate classifies 3,630 modules (3,597 trusted, 29 pinned by hash)
  in 12.6\,s. Separately, \code{lake exe axiomaudit} checks every one of the 1,501 declarations of
  the 31 organiser modules (contract, library, controls, bridge) against the three axioms, so that no
  organiser lemma a submission imports smuggles in \lean{sorry} or native evaluation.
\item \textbf{Admission and scoring.} Admission caps (image size, tape size, key length, message and
  index-oracle lengths, leaves, index-challenge widths, query count, round count, blow-up) are
  \emph{clauses of the certificate}, so the kernel has already checked them. So are the scores:
  the certificate proves that the rendered \lean{zk4} equals \lean{scoreZk4} of the rendered
  summaries and counters (T2: a rendered table of $S^4$ values; latency: \lean{scoreLat3}). The
  scorer only takes roots ($\mathrm{Score}_{\mathrm{zk}}=\mathit{zk4}^{1/48}$,
  \cref{sec:cost}). Because $c_k$ bounds each counter separately, a re-weighting of the profiles
  leaves the cost proof valid; only the one-line \lean{score} clause is re-checked
  (\cref{subsec:arch-versioning}).
\item \textbf{Badge path} (stage~3, \cref{subsec:arch-badge}), in its own sandboxed process in the
  gate's tree, after the gate. v0.3 provides the hook; the checks are planned.
\end{enumerate}

\begin{decision}{A-3}{The harness order is a rule}
\begin{enumerate}[nosep,leftmargin=1.3em]
\item Nothing built from the solution runs before comparator in any tree that comparator reads.
  Comparator runs first, in a fresh tree, with the organiser package, the packages and the toolchain
  outside the tree's writable \code{.lake}, read-only to the sandbox.
\item The gate runs afterwards, in a second fresh tree, and only after the manifest check.
\item The badge path runs last, as its own process; never inside the gate's process or an
  emulator's process.
\item The organiser \code{.olean} manifest is checked before the first stage and after every stage
  that runs solution code.
\item The gate and the badge operate on the build that comparator verified: every compiled
  artefact of the gate tree must be byte-identical to the comparator tree's (A7-M3).
\end{enumerate}
Any one of order, read-only mounts and the manifest rejects the known attack; all three are
required.
\end{decision}

\paragraph{The incident behind Decision~A-3 (A6-G0).} In v0.3 as first built, \code{ci.sh} built the
toy challenge \emph{and} solution in the organiser tree (the gate needs the solution's
\code{.olean}s), ran the gate, and only then ran comparator. Review A6 built a two-file mini
project with the real comparator binary (\path{probes/a6-probes/cmpproj/}): the stand-in contract
defines an unsatisfiable \lean{Good s := s.vk = 42 ∧ s.vk = 43}, and the solution imports a module
whose elaboration (an \lean{\#eval} running \lean{IO.Process}) atomically replaces
\path{.lake/build/lib/lean/RecursionFast/Contract.olean} by one with \lean{Good := True}. Lake did
not notice: its traces are computed from sources, so the next \code{lake build} rebuilt nothing,
and the challenge, elaborated for the first time against the tampered \code{.olean}, agreed with the
solution. Comparator printed \enquote{Your solution is okay!} for a false statement. In a fresh tree
with comparator first, the same solution is rejected (\enquote{Const does not match}). Comparator's
documentation states the precondition (\enquote{you have not previously tried to compile the
Solution file}), and a real landrun would not have helped, because comparator itself grants the
build write access to the project's \code{.lake}, where the contract \code{.olean}s lived. The lesson
holds for any Lean-checked competition: \emph{Lake does not detect a rewritten dependency
\code{.olean}}, so the trusted statement must be compiled where untrusted code has never run and
cannot write, and the trusted build must be pinned by content. The regression
\code{tests/pipeline/cmpproj.sh} (45\,s, in CI) keeps four variants: the old order still accepts (a
live reference); under the harness the attack is rejected by the stage-0 manifest check, by the
read-only organiser package (\enquote{read-only file system}, for A6's relative-path rewrite and for a
variant that rewrites the contract \code{.olean} wherever \code{LEAN\_PATH} finds it), and, without a
sandbox, by order alone.

\subsection{The badge path}\label{subsec:arch-badge}

The badge path is the only place where candidate code runs, and it gates records, never the
proved score. Its checks are specified here and not yet implemented (v0.3 has only the stage-3
hook). It runs after comparator and the gate, as its own sandboxed process (Decision~A-3):
a compiled indexer or prover is never loaded into the gate or the trusted emulator, so an
\lean{init} or \lean{export} entry that slipped past the gate could not act inside a trusted process.
In a sandbox with stated memory and time limits the harness (i)~evaluates the
\emph{compiled Lean indexer} on the hidden benchmark indices and requires the native indexer's
output to match byte for byte; (ii)~runs the native prover against harness-sampled challenges;
(iii)~on instances up to $2^{16}$ rows, compares the native prover's transcript with the compiled
Lean prover's; (iv)~runs the trusted emulator on the accepting transcripts and checks
$\textit{measured} \le \textit{claimed}$ componentwise; and (v)~enforces prover budgets of 1\,h
at $2^{20}$ and 8\,h at $2^{24}$ wall-clock, with a stated memory limit. Step~(i) closes a backdoor
found in review (A1-M3): without it, a verifier may accept when an index word equals a magic
constant that the Lean indexer never writes, and a native indexer that writes it passes every
demonstration with no work. A failure of~(iv) cannot happen if the theorem and the trusted
semantics are right; it indicates a bug in the machine semantics, the renderer or the emulator, and
halts the competition until resolved.

\subsection{What agents control and what is trusted}\label{subsec:arch-trusted}

A submitter controls every definition and proof in its root, the images and the claim. It does not
control the statement, which is the organizers' challenge, nor the relation, the fields, the
machine, the round shapes, the caps, the cost profiles and score functions, the benchmark or the
Fiat--Shamir prefix.
\Cref{tab:arch-trusted} lists everything a reader must accept.

\begingroup\small
\renewcommand{\arraystretch}{1.12}
\begin{longtable}{@{}>{\raggedright\arraybackslash}p{0.2\textwidth}>{\raggedright\arraybackslash}p{0.43\textwidth}>{\raggedright\arraybackslash}p{0.31\textwidth}@{}}
\caption{The trusted surface of recursion.fast.}\label{tab:arch-trusted}\\
\toprule
Item & What it is & Why trusted; mitigation \\
\midrule\endfirsthead
\toprule
Item & What it is & Why trusted; mitigation \\
\midrule\endhead
Contract files & 12 modules and a root, \textbf{796 code lines} (v0.2: 582), 471 comment lines, every definition with a docstring: fields $K$, $E$ and encodings; sparse R1CS, classes and padding maps; round shapes and transcripts; local security definitions; holographic \lean{Protocol}; pipeline timer; machine (decoder arms, W-op and Zbc semantics, system calls, harts, 18-counter first-touch accounting); \lean{Agrees}; cost clauses, zkVM profiles, score functions and admission caps; \lean{QProg}; accumulation; certificates. Listed in \cref{sec:appendix-contract} & statements of intent; protected-file pin; two-reviewer changes; controls must pass; imports Lean core, Mathlib and the vendored model only (no ArkLib, CompPoly, VCVio) \\
Challenge stubs & \code{T0}, \code{T1}, \code{T2} and \code{L1} (the latency-board stub, for \lean{Latency1}), with typed placeholders; everything except \lean{submission} is organiser data & reviewed with the contract; every stub rendered and elaborated in CI \\
Lean toolchain and checking stack & Lean v4.34.1 kernel; elaborator and frontend (the meaning of the statement is the organiser-elaborated challenge); Lake (resolves which \code{.olean} is loaded); compiler and runtime (gate, axiom audit, badge executables); comparator and \code{lean4export}; nanoda as second kernel; axiom whitelist; \code{leanchecker} (CI and submitter check only) & kernels have had bugs (three proof-of-\lean{False} routes fixed in v4.34.0 \cite{lean4340}); two independent kernels; re-verification on every toolchain bump; the comparator, \code{lean4export} and nanoda binaries are not yet hash-pinned \\
Imported definitions & Mathlib's \lean{Finset} (card, \lean{Ico}, \lean{image}), \lean{Set}, \lean{Fin} tuples (\lean{Fin.cons}, \lean{Fin.append}, \lean{Fin.snoc}, \lean{Fin.induction}, \lean{![…]}), \lean{List.ofFn}, \lean{Nat.card}, \lean{NNReal}, \lean{ℚ} with ceiling, \lean{List.Forall₂}, core \lean{BitVec}; \lean{Matrix} only in the untrusted dense bridge & stable, widely audited; comparator compares every used constant \\
Vendored RISC-V model & riscv-zkvm \code{Rv64/*} at \texttt{6810fc1}, byte-identical with a check script \cite{riscvzkvm2026}: \lean{MachineState}, \lean{execInstrBr}, base \lean{decode}; 474 reached declarations, 930 source lines (ZisK accelerator code dead at run time) & a hand model, not proved against the Sail specification; differential tests against QEMU for register ALU, M, Zbc, W-op and load/store semantics (control flow and fetch not covered) \\
Harness & \emph{implemented:} renderer and claim schema with path confinement; the stage order of \code{harness.sh} (Decision~A-3), the read-only organiser package and the \code{.olean} hash manifest; environment gate and axiom audit tool (Lean core only). \emph{Planned:} preflight and source policy; tape loader and oracle construction; trusted emulator with the 18 counters and the timer; canonical form of benchmark indices; Fiat--Shamir prefix and challenge sampler; benchmark generators and hidden seeds; badge checks. The scorer only takes roots of certified values & short, reviewed scripts with test vectors and attack probes; emulator to be differential-tested against the Lean semantics; sampler exact over binary fields \\
Host tools & Python 3.11+ (renderer, manifest, harness glue), bubblewrap (the v0.3 sandbox shim), \code{sha256sum}, the harness shell scripts & to be pinned in the production container (not yet pinned); short and reviewed \\
Cited meta-theorems & state-restoration and BCS/COS compilation: soundness by \cite[Thms.~31.2.1, 32.8.9]{CY24-Book}; knowledge by \cite[Thm.~B.4]{BCFW25-WARP} with \cite[Thm.~32.8.10]{CY24-Book} & published proofs, not formalized (ArkLib has no BCS transform); stated with side conditions (\cref{sec:security}) \\
Named hypotheses & conditional track only (\code{Hypotheses.lean}, planned; not in v0.3): versioned, uninhabited structures transcribing published theorems with regime guards, first \lean{McaJohnson} & two-reviewer admission; records flagged with the hypothesis version (\cref{subsec:sec-hypotheses}) \\
Deployment assumptions & the deployed hash behaves as a random oracle; \lean{RecursionHeuristic} for recursive deployment; classical adversaries; no zero knowledge & prose, never hypotheses of theorems \\
\bottomrule
\end{longtable}
\endgroup

\paragraph{Known gaps.} The following are accepted in v1 and tracked in \cref{sec:issues}.
\begin{itemize}[nosep]
\item The decoder and the Zbc semantics are trusted and tested, not proved against Sail.
\item There is no organiser emulator yet. The only one is the baseline's unverified C emulator,
  differential-tested against the Lean \lean{execute} for the 18 counters, first touch and the timer
  on small instances, not for harts; the trusted emulator must be built and tested the same way.
\item The v0.3 pipeline runs under a bubblewrap shim with landrun's command line, not under real
  landrun. The shim (\code{bwrap-landrun.sh}) mounts \code{/} read-only, makes only the fresh tree's
  \code{.lake} writable, and has no network and its own PID namespace; landrun's executability restrictions are not
  modelled. The packages directory is a shared, read-only copy of the frozen v0.2 tree's. Production
  must use the real sandbox with the same stage order.
\item The manifest covers \code{.olean} files, not native artefacts (\code{.c}, \code{.o},
  \code{.a}) of organiser modules. A badge executable must be linked against the read-only organiser
  package only. The manifest must be regenerated from a clean organiser build after every legitimate
  change of contract, vendored code, packages or toolchain.
\item The trust story depends on harness discipline: the order of Decision~A-3, the read-only
  mounts and the manifest are harness code, not Lean statements, and every review round so far has
  found a new hole at this level (\cref{subsec:thr-evidence}).
\item Declaration options such as \lean{debug.skipKernelTC} are invisible to the gate and are caught
  only by the kernel replays; two replays (comparator's Lean kernel and nanoda) are therefore
  mandatory. \code{leanchecker} replays in Lean's own kernel, so it adds a replay, not kernel
  diversity.
\item The gate assumes that compiled code equals the kernel definition unless an attribute says
  otherwise; other environment extensions carrying compiled code are not compared. A probe should
  confirm this before launch (A7-m17).
\item Benchmark indices are representations (entry lists); their canonical form is a harness
  convention outside the contract (\cref{subsec:proto-relation}).
\item The Fiat--Shamir prefix and the sampler are harness code. A bug there is a real attack with a
  valid certificate.
\item The compilation theorems are cited, not proved in Lean; at the base level the knowledge
  certificate is equivalent to round-by-round soundness (\cref{subsec:sec-extractor}); the
  accumulation track has one-step theorems only until the tree theorem exists (\cref{sec:recursion}).
\item Comparator's fresh-environment replay ignores declaration-local recursion limits. The
  contract therefore uses no local options, and its only computational certificate (CompPoly's
  Rabin certificate for the degree-64 modulus, in the untrusted library) replays unpatched in
  101\,s. Agents are warned that proofs relying on \lean{set\_option maxRecDepth} may fail replay.
\item Metaprograms may load modules at elaboration time. This is not restricted: soundness rests
  on export and replay, and confidentiality on the sandbox.
\item Comparator matches definition holes only by name and type. The certificate is therefore
  written so that \emph{any} term of the stated type is a legitimate win.
\end{itemize}

\subsection{Budgets}\label{subsec:arch-budgets}

Budgets are stated in resources, not lines: a single kernel-evaluated \lean{decide} table can cost
more than ten thousand ordinary lines, and single-threaded export and replay is a floor that more
cores do not lower.

\begin{table}[h]
\centering\small
\begin{tabularx}{\linewidth}{@{}lccX@{}}
\toprule
Track & Wall clock & CPUs & Further limits \\
\midrule
T0 (Lean only) & 30\,min & 8 & peak 32\,GiB \\
T1, T2 & 3\,h & 8 & elaboration $\le 16$ CPU-h; export + replay $\le 75$\,min single-threaded; peak RSS $\le 48$\,GiB aggregate and $\le 16$\,GiB per process \\
Badge path & per run & -- & prover 1\,h at $2^{20}$, 8\,h at $2^{24}$; stated memory limit \\
\bottomrule
\end{tabularx}
\caption{Check and badge budgets.}
\label{tab:arch-budgets}
\end{table}

\noindent Calibration: sig.golf's 154k-line proof needs 2.8 CPU-h of elaboration plus 23\,min of
single-threaded replay; at that rate 200k lines fit the T1 budget with about $2\times$ margin.
The certified baseline (1{,}920 lines, written without a library) passes one harness run in about
6\,min and the full check in 13.6\,min (\cref{subsec:base-time}); A3 had estimated 8\,min warm and
20--25\,min cold. Kernel-evaluated lookups and size checks against a 1\,MiB image peak at about
9\,GB (\cref{subsec:arch-binding}), so such proofs are kept in modules of their own, and Lake runs
with \code{-j8} under the per-process memory limit. Importing the contract takes
5.0\,s and 1.9\,GB; a clean build of the contract (including the vendored model) takes 80\,s.

\subsection{Versioning and record revalidation}\label{subsec:arch-versioning}

The \emph{contract id} is the SHA-256 digest of the file that pins the digests of all protected
files (contract modules, stubs, toolchain manifest, comparator configuration). Every verdict
records the contract id and the commit of the checking software. A contract change produces a new
id. Records on the old id stay on the board under their id until a maintainer ports them to the
new contract. Ported records are listed with their original and re-check provenance, and porting
does not transfer credit. Three events force revalidation of every affected record: (i)~a
toolchain or dependency bump; (ii)~a change of the machine semantics or the decoder; (iii)~the
removal or amendment of a named hypothesis, after which dependent conditional records are
re-checked or withdrawn. A hypothesis that becomes a Lean theorem is inhabited in the library;
the dependent records become unconditional on re-check. Re-weighting the cost vector, by contrast,
changes only the \lean{profiles} constant of \code{Cost.lean} (a new contract id): the expensive cost proofs stay valid,
and each record's one-line \lean{score} clause is re-proved against the re-rendered score
(\lean{norm\_num} or \lean{decide} on rationals).
